\chapter{More lattices}

After reading the LLL chapter, the reader might have possibly come to the conclusion that Regev collects a single sample - the first vector from LLL - from every run of the algorithm, and hopes that this vector happens to correspond to a non-trivial root. This, however, is far from the case. Certain procedures we have described up to now, like the gaussian superposition and the chance of hitting a trivial root rely on probabilistic outcomes, and we also need to account for the noise inherent in quantum systems. Regev therefore attempts to generate as many vectors as possible from a single run, to increase his probability of finding non-trivial roots.

Now, suppose we have our $d+k$ normalized d-dimensional samples of $\mathit \Lambda^*$. By normalized, we mean that we have divided the original samples by $D$ and now every entry is in $[0,1)$. This means we've mapped every sample to a unique point on the unit torus. Suppose now that each sample $\mathbf{w_1,w_2,\dots,w_{d+k}}$ differs from a real $\Lambda^*$ vector $\mathbf{v_1,v_2,\dots,v_{d+k}}$ by some quantity $\delta_i$. Let $\delta$ be some upper bound on the magnitude (absolute value) of the $\delta_i$, so $\delta_i \leq \delta$ and define the $2d+k$ \footnote{$k \geq 4$ from Lemma 6.3} dimensional lattice $\mathit \Lambda'$ spanned by the columns of $M$, where $M$ is the following matrix: 
\begin{equation}
M = \left[
\begin{array}{c|c}
  I_{d\times d} & 0 \\
  \hline
  \begin{matrix} \delta^{-1}\mathbf{w_1^T} \\  \vdots \\ \delta^{-1}\mathbf{w_{d+k}^T} \end{matrix} & \delta^{-1} \cdot I_{(d+k)\times (d+k)}
\end{array}
\right]    
\label{eq10.1}
\end{equation}



Regev proves using the Cauchy-Schwarz inequality the following lemma (combination of Lemma 4.4 and Corollary 4.5 in \cite{Regev2025}). \\

\textbf{Lemma 10.1}: For any $\mathbf{u} \in \Lambda$, there exists a vector $\mathbf{u}' \in \Lambda'$ whose first $d$ coordinates are equal to \textbf{u} with norm $\lVert \mathbf{u'} \rVert \leq \lVert \mathbf{u} \rVert \cdot \sqrt{1+d+k}$. Moreover, for $k = 4$, it holds that the first $d$ coordinates of every nonzero column vector $\mathbf{u'} \in \Lambda'$ with norm 
\begin{equation}
\lVert \mathbf{u'} \rVert < \frac{\delta^{-1}}{6 \cdot (4det(\Lambda))^{1/(d+k)}}  
\label{eq10.2}
\end{equation}
form a vector in $\Lambda$ with probability at least 1/4.

To understand why we are doing this, suppose we multiply some arbitrary $\mathbf{u} \in \mathbb{Z}^d$ with the first d columns of M. This results in a $2d+k$ dimensional vector $\mathbf{u'}$ whose first $d$ coordinates are the first $d$ coordinates of $\mathbf{u}$ and the next $d+k$ coordinates are the inner products $\langle \mathbf{u,w_i} \rangle$ scaled by $\delta^{-1}$. If $\mathbf{u} \in \Lambda$, those inner products are mathematically guaranteed to be really close to some integer multiplied by $\delta^{-1}$ - specifically within a distance of $\delta \rVert \mathbf{u} \lVert$. Now, suppose we define a vector $\mathbf{\hat{u}}$ that is equal to subtracting the closest integer to $\langle \mathbf{u,w_m} \rangle$ from each of the last $m$ coordinates of $u'$. 

Then the last $m$ coordinates of $\mathbf{\hat{u}}$ are either very small if $\langle \mathbf{u,w_m} \rangle = 0$, which corresponds to $u$ being in the dual lattice of $\Lambda^*$ (e.e $\Lambda$) or very large otherwise. Specifically, after the multiplication with $\delta^{-1}$, in the first case the size of each of the last $m$ coordinates is bounded above by $\rVert \mathbf{u} \lVert$, while in the second case it is bounded below by $\delta^{-1}\epsilon/2$, where $\epsilon = (4det(\Lambda)^{-1/(d+k)})/3$ \eqref{eq10.2}. All this means is that this transformation makes the dual lattice vectors small and the non dual lattice vectors large.

Since any vector $\mathbf{u'} \in \Lambda'$ is a linear combination of the vectors formed by the columns of $M$, finding the shortest vectors in $\Lambda'$ naturally produces our desired solution. For valid vectors $\mathbf{u} \in \Lambda$, the last $m$ coordinates of $\mathbf{u'}$ will evaluate to very small values, meaning we are just isolating the shortest $\mathbf{u}$.The process of finding short vectors automatically filters out $\mathbf{u} \notin \Lambda$ by virtue of their last $m$ coordinates having really large values. This means that running LLL on $M$ will provide us with at least one vector whose first $d$ coordinates correspond to a short vector in $\Lambda$.

For now, from this correspondence, we can see how from our noisy samples of $ \mathit \Lambda^*$ we will produce a clean, short base for $\mathit \Lambda$ sufficient to give us a non-trivial factor. Suppose we run LLL on the basis formed by M, returning some vectors $\mathbf{b_1, b_2, \dots, b_{d+k}}$ and then Gram-Schmidt orthogonalization on the result, giving us some orthogonal vectors $\mathbf{b^*_1, b^*_2, \dots, b^*_{d+k}}$. Let $T$ be some norm bound, and let $h \leq k$ be the largest number such that 
\begin{equation}
    \lVert\mathbf{b^*_h} \rVert <  2^{k/2} T
\end{equation}

Then, according to Claim 5.1 in Regev's paper, any vector in the sublattice $\mathit \Lambda'_T \subseteq \mathit \Lambda'$ formed by vectors of norm $\leq T$ can be generated by $\mathbf{b_1}, \mathbf{b_2}, \dots ,\mathbf{b_h}$. 
Due to \textbf{Lemma 10.1}, the first d coordinates of the $\mathbf{b_i}$ correspond to vectors in $\mathit \Lambda$. This gives us not just one, but multiple short vectors in $\mathit \Lambda$ that could also be in $\mathit \Lambda / \mathit \Lambda_0$. In fact, at least one of those vectors is guaranteed to not be in $\mathit \Lambda_0$, otherwise $\mathit \Lambda_T$  would be equal to $\mathit \Lambda_0$ itself (since $\mathit \Lambda_0$ is a group with closure). Therefore, taking $k \geq 4$ (to overshoot the >1/4 probability that each one of them is in $\mathit \Lambda_0$), we can reliably end up with a vector corresponding to a non-trivial factor.

Lastly, note that the reason we normalized our samples is that it really doesn't matter how large the original sample is. It is very easy to prove that the inner product of any arbitrary vector with a normalized vector of the dual is still an integer, and thus our previous logic still holds.











